\documentclass[10pt,a4paper]{amsart}
\usepackage[margin=19mm]{geometry}
\usepackage{amsmath,amssymb,mathtools}
\usepackage[scr=rsfs,cal=euler]{mathalfa}
\usepackage{xfrac}
\usepackage[unicode,hidelinks,bookmarks=false]{hyperref}
\newcommand{\ie}{i.e.\ }
\newcommand{\cf}{cf.\ }
\newcommand{\resp}{respectively\ }
\newcommand{\Subsection}{\subsection}
\providecommand{\nobreakdash}{\nobreak\hbox{-}}
\setlength{\emergencystretch}{2em}
\multlinegap0pt
\usepackage{mathtools}
\usepackage{amssymb}
\usepackage{float}
\usepackage{csquotes}
\newtheorem{lemma}{Lemma}
\newtheorem{proposition}{Proposition}
\usepackage{cleveref}
\crefname{lemma}{Lemma}{Lemmas}
\crefname{proposition}{Proposition}{Propositions}
\newcommand{\F}{\mathbb{F}}
\newcommand{\GL}{\mathrm{GL}}
\newcommand{\M}{\mathrm{M}}
\newcommand{\Z}{\mathbb{Z}}
\newcommand{\diag}{\mathrm{diag}}
\title{Periodic Points of Power Maps in Finite Matrix Groups and Algebras}
\author{Saikat Panja}
\date{Source: arXiv:2603.12295v1 (CC BY 4.0)}
\begin{document}
\maketitle

\noindent\emph{Source and licence.} Periodic Points of Power Maps in Finite Matrix Groups and Algebras. Version arXiv:2603.12295v1 (CC BY 4.0), distributed by arXiv under the Creative Commons Attribution 4.0 International licence, http://creativecommons.org/licenses/by/4.0/, as recorded in the arXiv licence metadata for this exact version and shown on the abstract page https://arxiv.org/abs/2603.12295v1. Full licence notice: \texttt{licenses/arxiv-2603.12295-CC-BY-4.0.txt}.\par\smallskip

\noindent\textit{Editorial scope.} Counted unit: the article's proposition prop:candidacy (source0.tex lines 260--264). The article's complete original proof of it is delivered with it (source0.tex lines 265--327, one \texttt{proof} environment of 63 lines). No mathematics is written, repaired or re-derived: the statement and the proof are the article's own lines, in the article's own notation, and no retained line is substituted or re-flowed.  Retained uncounted as the source-local context the proof needs: lines 250--252.  Nothing was omitted from any retained span, so the package carries no omission record.  No retained line contains a citation, so the wrapper carries no bibliography and no bibliography entry.  No retained line contains a figure environment, an image-inclusion command or drawing code, so no image is delivered and the package declares no asset dependency.  Editorial formatting only, in addition: an AMS \texttt{amsart} preamble that restates the article's own declarations and macros used by the retained lines, namely the environments \texttt{lemma}, \texttt{proposition} on separate counters, together with the standard packages those lines need.  The retained environments print in this document as Lemma 1, Proposition 1 in order of appearance, whereas the article prints the same items with the numbering of its own full document.  The complete native source of the article as distributed in the arXiv e-print for this exact version is bundled unchanged as \texttt{sources/196269/source0.tex}.
\begin{lemma}\label{lem:conj}
    Let $f\in \F_q[t]$ be a polynomial, $A\in\M_n(q)$ be a matrix. Then $f^m(A)=A$ iff for any conjugate $B$ of $A$ one has $f^m(B)=B$.
\end{lemma}

\begin{proposition}\label{prop:candidacy}
Let $L$ be a prime integer coprime to the characteristic $p$ of the finite field $\F_q$, and let $f$ denote the map $x\mapsto x^L$. Assume that $\delta_L(p)\mid d$, where $q=p^d$ for some $d\ge 1$.
A matrix $A \in \M_n(q)$ is periodic if and only if all eigenvalues of $A$ are periodic and there is no nilpotent block in the Jordan canonical form of $A$, over an algebraic closure of $\F_q$.
In particular, a matrix $A\in \GL_n(q)$ is periodic if and only if all its eigenvalues are periodic.
\end{proposition}

\begin{proof}
    We work in the field $\overline{\F_q}$, invoking \Cref{lem:conj}. 
    Hence, the matrix $A$ can be taken into its Jordan canonical form $\diag(J_{\lambda_q,m_1},J_{\lambda_2,m_2},\cdots, J_{\lambda_u,m_u})$ for positive integers $m_i$, where
    \begin{align*}
        J_{\lambda,m}=\begin{pmatrix}
            \lambda & 1 &&&&\\
            &\lambda&1&&&\\
            &&\ddots&&&\\
            &&&\lambda&1\\
            &&&&\lambda
        \end{pmatrix}_{m\times m}\in\M_m(\overline{\F_q}).
    \end{align*}
    Since $f^m(\diag(J_{\lambda_q,m_1},J_{\lambda_2,m_2},\cdots, J_{\lambda_u,m_u}))=\diag(J_{\lambda_q,m_1},J_{\lambda_2,m_2},\cdots, J_{\lambda_u,m_u})$ if and only if $f^m(J_{\lambda_i,m_i})=J_{\lambda_i,m_i}$ for all $i$, we may well assume that $A$ up to conjugacy has a single Jordan block.
    When $\lambda=0$, there is nothing to prove as $f^r(J_{0,n})=0$ for some $r>0$; hence, such a matrix can never enter a cycle. 
    So we assume that $\lambda\neq0$ and consequently $J_{\lambda,n}\in\GL_n(q)$.
    Now, note that for any $r>0$
    \begin{align*}
        J_{\lambda,n}^r=
\begin{pmatrix}
\lambda^r & \dbinom{r}{1}\lambda^{r-1} & \dbinom{r}{2}\lambda^{r-2} & \cdots & \dbinom{r}{n-1}\lambda^{r-(n-1)} \\
0 & \lambda^r & \dbinom{r}{1}\lambda^{r-1} & \cdots & \dbinom{r}{n-2}\lambda^{r-(n-2)} \\
0 & 0 & \lambda^r & \cdots & \dbinom{r}{n-3}\lambda^{r-(n-3)} \\
\vdots & \vdots & \vdots & \ddots & \vdots \\
0 & 0 & 0 & \cdots & \lambda^r
\end{pmatrix},
\end{align*}
which implies that for having $J_{\lambda,n}^r=J_{\lambda,n}$ 
\begin{align*}
    \displaystyle{\dbinom{r}{i}}\lambda^{r-i}\equiv\begin{cases}
       \lambda&\text{when }\,i=0\\
       1&\text{when }\,i=1\\
       0&\text{when }\, n-1\geq i\geq 2
    \end{cases}.
\end{align*}
Thus, we get that $p\bigm|\dbinom{r}{i}$ for all $i\geq 2$, and $r\equiv 1\pmod{p}$.
If $r=L^{r'}$ we have that $\delta_p(L)|r'$.

Consider the base-$p$ expansion, $r=1+\sum\limits_{i=1}^{u}r_ip^i$ and $s=\sum\limits_{i=0}^{u}s_ip^i$.
Recalling Lucas theorem, one has
\begin{align*}
    \dbinom{r}{s}\equiv\prod\limits_{i=1}^u\dbinom{r_i}{s_i}\pmod{p}.
\end{align*}
First we get that $r\geq n$; otherwise $\dbinom{r}{n-t}=1$ for some $t\geq 1$.
Write $r=p^{v_p(r-1)}a$ for some $a\not\equiv 0 \pmod{p}$.
Let $a=\sum\limits_{i=0}^\ell a_ip^i$.
Thus we have 
\begin{align*}
    r=1+a_1p^{v_p(r-1)+1}+a_2p^{v_p(r-1)+2}+\ldots=1+\sum\limits_{i=0}^{\ell}a_ip^{v_p(r-1)+i};
\end{align*}
where $v_p(r-1)\geq 1$ as $p|r-1$, and $a_1\neq 0$. Now, if possible, let $n>p^{v_p(r-1)}$. Fix $i_0=1+p^{v_p(r-1)}$, which implies by Lucas theorem that
\begin{align*}
    \dbinom{r}{i_0}\equiv\dbinom{1}{1}\cdot\dbinom{a_1}{1}\neq 0\pmod{p};
\end{align*}
also $2<i_0\leq n-1$. 
We conclude that $n\leq p^{v_p(r-1)}$.
By applying Lucas theorem again, when $n\leq p^{v_p(r-1)}$, one has $\dbinom{r}{i}\equiv0\pmod{p}$ for all $2\leq i\leq n-1$.

Now, let us focus on the power map $x\mapsto x^L$, and we work over the algebraic closure of the underlying field.
A matrix $A=\bigoplus\limits_{i=1}^{\ell}J_{\lambda_i,m_i}$ is periodic- that is $A^{L^r}=A$ if and only if $\lambda_j^{L^r}=\lambda_j$ for all $1\leq j\leq \ell$ and $m_i\leq p^{v_p(L^r-1)}$.
One also has that, if $\mu^{L^r}=\mu$ for some $\mu\in\overline{\F_q}$, then for any $b\in\Z_{\geq 1}$, $\mu^{L^{br}}=\mu$. 
Using the Lifting The Exponent, one has $v_p(L^{br}-1)=v_p(L^r-1)+v_p(b)$. 
Given $n$, one may choose $b$ such that $n<v_p(L^{br}-1)$, where $\lambda^{L^r}=1$ with $\lambda\in\sigma(A)$. The result follows.
\end{proof}

\end{document}
